\begin{table}
    \footnotesize
    \centering
    \begin{threeparttable}
        \caption{LLM readability of authors' \(t\)th paper (draft and final)}
        \label{table9llm}
        \begin{tabular}{p{3cm}S@{}S@{}S@{}S@{}S@{}S@{}}
            \toprule
            &{\(t=1\)}&{\(t=2\)}&{\(t=3\)}&{\(t=4\text{--}5\)}&{\(t\ge6\)}\\
            \midrule
            \multicolumn{6}{l}{\textbf{Predicted \(R_{jP}-R_{jW}\)}}\\\quad Women                   &        0.42***&        0.39***&        0.36***&        0.31***&        0.33** \\
                                          &      (0.10)   &      (0.09)   &      (0.09)   &      (0.11)   &      (0.14)   \\
            \quad Men                     &        0.42***&        0.43***&        0.43***&        0.41***&        0.47***\\
                                          &      (0.02)   &      (0.01)   &      (0.01)   &      (0.01)   &      (0.02)   \\
            \midrule\multicolumn{6}{l}{\textbf{Marginal effect of female ratio}}\\\quad Published article       &        0.43***&        0.44***&        0.45***&        0.46***&        0.46***\\
                                          &      (0.11)   &      (0.08)   &      (0.07)   &      (0.08)   &      (0.10)   \\
            \quad Draft paper             &        0.44***&        0.48***&        0.52***&        0.56***&        0.60***\\
                                          &      (0.14)   &      (0.11)   &      (0.10)   &      (0.11)   &      (0.14)   \\
            \midrule\textbf{Diff.-in-diff.}&       -0.01   &       -0.04   &       -0.07   &       -0.10   &       -0.14   \\
                                          &      (0.11)   &      (0.10)   &      (0.10)   &      (0.12)   &      (0.14)   \\
            \bottomrule
        \end{tabular}
        \begin{tablenotes}
            \tiny
            \item \textit{Notes}. Sample 4,289 observations. Panel one displays magnitude of predicted \(R_{jP}-R_{jW}\) (the direct effect of peer review) for women and men over increasing \(t\). Panel two estimates the marginal effect of an article's female ratio (\(\beta_1+\beta_2\times t\)), separately for draft papers and published articles. Figures from FGLS estimation of~equation~(15), weighted by \(N_{it}\) (see~the data appendix). Control variables include citation count (asinh), \(\text{max. }T\) (author prominence) and \(\text{max. }t\) (author seniority), native speaker and editor and journal-year fixed effects. Standard errors clustered by editor and robust to cross-model correlation in parentheses. ***, ** and * statistically significant at 1\%, 5\% and 10\%, respectively.
        \end{tablenotes}
    \end{threeparttable}
\end{table}